\documentclass[11pt,a4paper]{article}

% Cryptographic Hash Functions Laboratory - student report (pdfLaTeX, Overleaf / Tectonic).
\usepackage[T1]{fontenc}
\usepackage[utf8]{inputenc}
\usepackage{lmodern}
\usepackage[margin=2.1cm]{geometry}
\usepackage{microtype}
\usepackage{parskip}
\usepackage{amsmath,amssymb}
\usepackage{booktabs}
\usepackage{tabularx}
\usepackage{array}
\usepackage{float}
\usepackage{needspace}
\usepackage{enumitem}
\usepackage{xcolor}
\usepackage{fancyvrb}
\usepackage{cite}
\renewcommand{\citepunct}{], [}
\renewcommand{\citedash}{]--[}
\usepackage{hyperref}
\hypersetup{colorlinks=true, linkcolor=blue!50!black, urlcolor=blue!60!black, citecolor=blue!50!black,
  pdftitle={Cryptographic Hash Functions in Practice}, pdfauthor={Diellart Mulolli}}
\setlist{itemsep=0.2em, topsep=0.3em}
\setlength{\emergencystretch}{2em}
\newcolumntype{L}{>{\raggedright\arraybackslash}X}
\newcolumntype{P}[1]{>{\raggedright\arraybackslash}p{#1}}

\title{\textbf{Cryptographic Hash Functions in Practice}\\[0.3em]
\large From the birthday bound to HMAC, password storage, and real collisions}
\author{Diellart Mulolli (230756200035)\\[0.2em] \small Code and LaTeX source: \url{https://github.com/Diellart-Mulolli/hash-lab}}
\date{\today}

\begin{document}
\maketitle

\section*{Introduction}
A cryptographic hash function compresses a message of any length into a short, fixed-length digest and
is the quiet workhorse of modern security: it protects stored passwords, fingerprints software
downloads, forms the core of HMAC and of every digital signature, and chains the blocks of a
blockchain. Its usefulness rests on three properties --- preimage, second-preimage and collision
resistance --- and on the fact that a well-designed $n$-bit hash costs about $2^n$ work to invert but
only $2^{n/2}$ to collide, the birthday bound. This report works through the mathematics on a toy
checksum, surveys the real status of MD5 through SHA-3, shows a length-extension forgery and HMAC by
hand, estimates password-cracking and Merkle-tree costs, analyses a documented real-world failure (the
2017 SHAttered SHA-1 collision), and finally chooses hash-based mechanisms for a long-lived university
system. Six small Python scripts were used to verify every calculation; they are published, with the
LaTeX source and a running guide, at \url{https://github.com/Diellart-Mulolli/hash-lab}. Appendix~A
lists the values chosen in each script and its exact output.

%=====================================================================
\section*{Foundations}

\subsection*{Question 1 --- Security properties, and why a checksum is not a hash}
\textbf{(a)} The three properties, in plain terms:
\begin{itemize}
\item \textbf{Preimage resistance (one-way).} Given a digest $y$, it is infeasible to find any $x$ with
$h(x)=y$. This protects \emph{stored passwords}: a server keeps $h(\text{password})$, and an attacker who
steals the database still cannot recover the passwords by inverting the hash.
\item \textbf{Second-preimage resistance.} Given a \emph{fixed} $x$, it is infeasible to find a different
$x'$ with $h(x')=h(x)$. This protects a \emph{published software fingerprint}: the vendor pins the hash of
one specific installer, and an attacker cannot craft a different file with the same hash.
\item \textbf{Collision resistance.} It is infeasible to find \emph{any} pair $x\neq x'$ with
$h(x)=h(x')$. This protects \emph{a CA-signed certificate or any digital signature}: the signer signs
$h(M)$, so if the attacker can find two messages with the same hash (one benign, one rogue), a signature
on the benign one is also valid on the rogue one.
\end{itemize}
\textbf{Collision resistance is hardest}, because the attacker may choose \emph{both} messages and so
enjoys the birthday effect: a collision costs only $\approx 2^{n/2}$ work, against $\approx 2^{n}$ for a
(second-)preimage, where the attacker is pinned to one target. This is exactly why collisions fell first
for MD5 and SHA-1 while preimages did not.

\textbf{(b)} Toy hash $h(X)=(3X_0+2X_1+X_2)\bmod 256$, verified by \texttt{q1\_toy\_hash.py}
(Appendix~A.1).
\begin{enumerate}[label=(\roman*)]
\item $\mathrm{BID}=(66,73,68)$, so $h(\mathrm{BID}) = (3\cdot 66 + 2\cdot 73 + 68)\bmod 256 = 412\bmod
256 = \boxed{156}$.
\item \emph{Second preimage:} add $1$ to $X_0$ (which adds $3$) and subtract $3$ from $X_2$. This gives
$(67,73,65)=\texttt{"CIA"}$ with $h=156$. (The script confirms $256^2=65536$ messages share this hash.)
\item \emph{Targeted preimage starting \texttt{OK}:} with $X_0=79,\ X_1=75$ fixed, solve for $X_2$:
$X_2 = (156 - 3\cdot 79 - 2\cdot 75)\bmod 256 = 25$. So $(79,75,25)$ hashes to $156$.
\item Both (ii) and (iii) are easy because the formula is \textbf{linear (affine) and invertible modulo
$256$}: each output bit depends smoothly on the input, there is no diffusion or one-wayness, and one
byte can always be solved for from a single linear equation. The short 8-bit output is not the point ---
even a long linear checksum would fall the same way.
\end{enumerate}

\textbf{(c)} With $P(k)\approx 1-e^{-k(k-1)/2^{\,n+1}}$ for $n=8$ (\texttt{q2\_birthday.py}, Appendix~A.2):
\[
P(16) = 1-e^{-16\cdot 15/2^{9}} = 1-e^{-240/512} \approx 0.374\ (37.4\%).
\]
The rule of thumb $k\approx 1.18\cdot 2^{n/2}=1.18\cdot 16 \approx 18.9$ gives $P\approx\tfrac12$; the exact
smallest $k$ with $P\ge 0.5$ from the formula is $20$. A hash output must be \textbf{about twice as long
as a symmetric key} for equal security because the best generic attack on the hash is the birthday
collision at $2^{n/2}$, whereas brute-forcing a $\kappa$-bit key costs $2^{\kappa}$; matching them needs
$n/2=\kappa$, i.e.\ $n=2\kappa$ (so 256-bit SHA-256 pairs with 128-bit AES).

\textbf{(d)} A CRC-32 is a \emph{linear} function over $\mathrm{GF}(2)$: $\mathrm{CRC}(a\oplus b) =
\mathrm{CRC}(a)\oplus\mathrm{CRC}(b)$ (up to fixed constants). A stream cipher also encrypts by XOR with a
keystream. So an attacker who flips chosen plaintext bits knows \emph{exactly} which CRC bits flip, and
can flip the corresponding bits of the encrypted CRC to keep it consistent --- all without the key and
without decrypting. The checksum therefore detects random noise but not an intelligent attacker; this is
the WEP integrity flaw (Stamp, Section~5.5). Integrity needs a \emph{keyed} MAC, not a linear checksum.

\subsection*{Question 2 --- Output length and the status of real hash functions}
\textbf{(a)} Output lengths and current status (collision figures from the cited cryptanalysis):
\begin{table}[H]
\centering\small
\begin{tabularx}{\textwidth}{@{}lccP{1.5cm}LP{2.0cm}@{}}
\toprule
Hash & Output & Construction & Generic collision & Best known collision attack (year, cost) & Status today \\
\midrule
MD5 & 128 & Merkle--Damg\aa rd & $2^{64}$ & identical-prefix in seconds ($\le 2^{24}$); chosen-prefix $\approx 2^{39}$ (2009) \cite{stevenscp} & Broken; forbidden \\
SHA-1 & 160 & Merkle--Damg\aa rd & $2^{80}$ & identical-prefix $2^{63.1}$ (2017, SHAttered) \cite{shattered}; chosen-prefix $2^{63.4}$ (2020) \cite{shambles} & Broken; retired by 2030 \cite{nistsha1} \\
SHA-256 & 256 & Merkle--Damg\aa rd & $2^{128}$ & 31 of 64 steps practical (2024) \cite{sha2coll}; full function unbroken & Secure \\
SHA-512 & 512 & Merkle--Damg\aa rd & $2^{256}$ & 31 of 80 steps, $2^{97.3}$ (2024) \cite{sha2coll}; full function unbroken & Secure \\
SHA3-256 & 256 & Sponge (Keccak) & $2^{128}$ & 5 of 24 rounds practical \cite{sha3coll}; full function unbroken & Secure \\
\bottomrule
\end{tabularx}
\caption{Hash status. ``Steps/rounds'' attacks are on round-reduced versions; the standard functions are intact.}
\end{table}

\textbf{(b)} At $10^{12}$ hashes/s (\texttt{q2\_birthday.py}):
$2^{64}$ takes $\approx 1.8\cdot 10^{7}$\,s $\approx \mathbf{214}$ \textbf{days};
$2^{80}$ takes $\approx \mathbf{3.8\cdot 10^{4}}$ \textbf{years};
$2^{128}$ takes $\approx \mathbf{1.1\cdot 10^{19}}$ \textbf{years} (far beyond the age of the universe).
The generic collision cost of SHA-1 is $2^{80}\approx 38{,}000$ years at this rate, yet the first real
SHA-1 collision (2017) cost only $\approx 2^{63.1}$ compression calls --- about 6{,}500 CPU-years plus
110 GPU-years, i.e.\ a large but affordable GPU budget. The difference is the whole point: cryptanalysis
cut the cost by a factor of $\approx 2^{80-63.1}\approx 2^{17}$ ($\sim$130{,}000$\times$), so the
``generic $2^{80}$'' figure stopped describing SHA-1 once the function was broken.

\textbf{(c)} In the birthday attack on signatures (Stamp, Section~5.4), Trudy prepares about $2^{n/2}$
innocuous variants of a ``good'' message and $2^{n/2}$ variants of an ``evil'' message; by the birthday
bound a \emph{good} variant and an \emph{evil} variant are likely to share a hash. She has Alice sign the
good variant, and the same signature is valid on the colliding evil one --- so she needs only $2^{n/2}$
variants of each, not $2^{n}$. A contract with 32 independent either/or phrasings yields $2^{32}$
variants, which defeats any hash with $n\le 64$ bits, since $2^{64/2}=2^{32}$. \emph{Link to last week:}
ECDSA signs $h(M)$, not $M$. P-256 gives $\approx 128$-bit security, so it is paired with SHA-256 whose
$2^{128}$ collision resistance matches it and is not the weak link. If MD5 were used instead, Trudy could
collide a benign and a malicious contract and transfer Alice's valid P-256 signature onto the malicious
one; the strong curve would not help, because the signature only ever commits to the hash.

\textbf{(d)} \emph{Grover's} algorithm searches an unstructured space quadratically faster, so it lowers
preimage resistance from $2^{n}$ to $\approx 2^{n/2}$ quantum evaluations. The \emph{Brassard--H\o
yer--Tapp (BHT)} algorithm lowers collision search from $2^{n/2}$ to $\approx 2^{n/3}$, but only with
$\approx 2^{n/3}$ quantum memory, which makes it uncompetitive with classical parallel collision search
in practice. SHA-256 and SHA3-256 stay adequate because $256$ bits leaves Grover preimage resistance at
$2^{128}$ and BHT collision resistance at $\approx 2^{85}$ (and higher in realistic models) --- still
$\ge 128$-bit post-quantum security. The hash-based post-quantum signature standards are
\textbf{SLH-DSA} (FIPS~205, the SPHINCS\textsuperscript{+} family) \cite{fips205} and the stateful
\textbf{LMS/XMSS} (NIST SP~800-208) \cite{sp800208}. Their security rests \emph{only} on the
(second-)preimage resistance of the underlying hash --- no number-theoretic assumption like factoring or
discrete logs --- so they remain safe as long as the hash does, which is exactly what makes them
conservative choices for the quantum era.

%=====================================================================
\clearpage
\section*{Constructions and Message Authentication}

\subsection*{Question 3 --- Inside the hash: Merkle--Damg\aa rd, the sponge, and length extension}
\textbf{(a)} In \textbf{Merkle--Damg\aa rd} (MD5, SHA-1, SHA-2) the message is padded (a $1$ bit, zeros,
and the 64-bit length) and split into blocks $B_1,B_2,\dots$; a compression function $F$ updates a
chaining value from a fixed constant $A$: $h = F(F(A,B_1),B_2)$ (Stamp, eq.~5.1). The output \emph{is the
entire final chaining value}, so the digest reveals the full internal state. In the \textbf{sponge}
(SHA-3/Keccak) a 1600-bit state splits into a rate part ($r=1088$ for SHA3-256) and a capacity part
($c=512$); each block is XORed into the rate and the permutation $f$ is applied (absorbing), then the
output is read from the rate (squeezing). The output is only part of the state: the \emph{capacity is
never output}. This difference is decisive for the next part.

\textbf{(b)} Length-extension forgery against $\mathrm{tag}=\mathrm{SHA\text{-}256}(K\|M)$, demonstrated
end-to-end in \texttt{q3\_length\_extension.py} (Appendix~A.3; a from-scratch SHA-256 that can resume from
a captured digest, cross-checked against \texttt{hashlib}).
\begin{enumerate}[label=(\roman*)]
\item \emph{What Mallory needs:} the message $M$, its tag, the desired suffix \texttt{\&to=trudy}, and the
\emph{length} of $K$ (16 bytes here). She does \textbf{not} need $K$. Because the digest equals the full
chaining value after $K\|M\|\text{pad}$, she can set that state and keep hashing.
\item \emph{Padding:} $K\|M$ is $16+18=34$ bytes. SHA-256 appends $\texttt{0x80}$, then zero bytes, then
the 64-bit big-endian bit length $34\cdot 8 = 272 = \texttt{0x110}$, which is $30$ bytes of ``glue''
padding, filling the first 64-byte block. This glue must appear \emph{inside} Mallory's forged message,
because the server computes $\mathrm{SHA\text{-}256}(K\|M')$ and only reproduces the captured state if
$M'$ literally contains those glue bytes that the original hash already absorbed.
\item \emph{Why the server accepts:} the server computes
$\mathrm{SHA\text{-}256}(K\|M\|\text{glue}\|\texttt{\&to=trudy})$. Since the tag is exactly the state
after the first block $K\|M\|\text{glue}$, resuming from the tag and absorbing the suffix yields the same
digest. The script's forged tag equals the server's recomputed tag
(\texttt{6abffc5f\ldots7723c1ac}), and because the server uses the \emph{last} \texttt{to=} value, the
request now pays \texttt{trudy}.
\item \emph{Unknown key length:} Mallory simply tries $|K|=0,1,2,\dots$; each guess is one candidate
forged request, and she knows she succeeded when the server accepts one. This is a tiny search (a few
dozen tries).
\end{enumerate}

\textbf{(c)} The attack fails for: \textbf{SHA3-256}, because the sponge never outputs the capacity, so
the digest does not reveal the state needed to resume; \textbf{SHA-512/256}, because it truncates its
output, so the 256-bit digest exposes only half the internal state and cannot be used as a resumable
chaining value; and \textbf{HMAC-SHA-256}, because the tag is $h((K\oplus opad)\,\|\,h(\cdots))$ --- an
outer keyed hash wraps the inner one, so the tag is not the resumable state of any key-prefixed message.
Stamp also warns (eq.~5.3) against $h(M\|K)$: there, if the attacker finds \emph{any} collision
$h(M)=h(M')$ (feasible for a broken hash, and found offline without the key), then
$h(M\|K)=h(M'\|K)$ for every $K$, so the construction inherits the hash's collision weakness directly.

\needspace{12\baselineskip}
\subsection*{Question 4 --- Hash, MAC, or signature: choosing and tracing integrity mechanisms}
\textbf{(a)}
\begin{table}[H]
\centering\small
\begin{tabularx}{\textwidth}{@{}P{4.2cm}LLL@{}}
\toprule
Property & Plain hash $h(M)$ & HMAC-SHA-256 & ECDSA P-256 / SHA-256 \\
\midrule
Key material & none & shared secret $K$ & private key (signer), public key (verifier) \\
Who can create it & anyone & holders of $K$ & holder of the private key only \\
Who can verify it & anyone & holders of $K$ & anyone, with the public key \\
Detects random errors? & yes & yes & yes \\
Detects a deliberate active attacker? & \textbf{no} (attacker recomputes $h$) & yes & yes \\
Authenticates the origin? & no & yes (to key holders) & yes (to one specific signer) \\
Non-repudiation? & no & no (both sides share $K$) & yes (only the signer has the key) \\
One real use & file dedup, download integrity vs.\ noise & TLS record MAC, API request auth & code signing, X.509 certificates \\
\bottomrule
\end{tabularx}
\end{table}

\textbf{(b)} HMAC-SHA-256 with $K=\texttt{key}=\texttt{0x6B\,0x65\,0x79}$, structure verified against the
standard library in \texttt{q4\_hmac.py} (Appendix~A.4), $B=64$, $ipad=\texttt{0x36}$, $opad=\texttt{0x5C}$.
\begin{enumerate}[label=(\roman*)]
\item $K$ is shorter than $B$, so it is right-padded with zero bytes to \textbf{64 bytes}.
\item $K\oplus ipad$: first three bytes $\texttt{0x5D\,0x53\,0x4F}$ (e.g.\ $\texttt{0x6B}\oplus\texttt{0x36}
=\texttt{0x5D}$), and the remaining 61 bytes are $\texttt{0x36}$ (since $\texttt{0x00}\oplus
ipad=ipad$). $K\oplus opad$: first three bytes $\texttt{0x37\,0x39\,0x25}$, remaining 61 bytes
$\texttt{0x5C}$.
\item $\mathrm{HMAC}(M,K)=\mathrm{SHA\text{-}256}\big((K\oplus opad)\,\|\,\mathrm{SHA\text{-}256}((K\oplus
ipad)\,\|\,M)\big)$.
\item SHA-256 adds at least 9 bytes of padding, so a 1000-byte message needs
$\lceil(1000+9)/64\rceil=16$ blocks. HMAC's inner hash covers $(K\oplus ipad)\,\|\,M$, one extra block,
$\lceil(1064+9)/64\rceil=17$ blocks; the outer hash covers $(K\oplus opad)$ plus the 32-byte inner
digest, $\lceil(96+9)/64\rceil=2$ blocks. Total $17+2=\mathbf{19}$ blocks versus $16$ for a plain hash:
about two extra compressions, \textbf{negligible} overhead.
\end{enumerate}

\textbf{(c)} A SHA-256 checksum on the download page protects only against \emph{accidental} corruption
--- truncated or bit-rotted downloads --- because anyone, including an attacker, can recompute $h$. It
does \textbf{not} protect against a deliberate attacker who replaces the installer \emph{and} the
checksum on the same page (a server compromise or a man-in-the-middle rewrites both; the plain-hash
column of (a) shows ``authenticates origin: no''). A better mechanism is a \textbf{digital signature}
(code signing with ECDSA/Ed25519 over SHA-256), with the vendor's public key distributed out-of-band in
the OS or browser trust store: only the vendor holds the private key, so the attacker cannot forge a
valid signature, and (a) shows this is the only column with origin authentication and non-repudiation
under a key the attacker does not have.

\textbf{(d)} HMAC appears in \textbf{HKDF} (TLS 1.3 key derivation). HKDF-Extract computes
$\mathrm{PRK}=\mathrm{HMAC}(\text{salt},\text{ECDH secret})$ to condense the non-uniform shared secret
into a uniformly random pseudorandom key, and HKDF-Expand chains HMAC calls to derive several independent
keys bound to context labels. HMAC's contribution is a \emph{keyed pseudorandom function} whose output
is computationally indistinguishable from random, which is what lets one Diffie--Hellman secret safely
produce many session keys. (In RFC~6979 the same HMAC-as-PRF property instead derives the ECDSA nonce
deterministically from the private key and $h(M)$, removing the need for a random number generator.)

%=====================================================================
\section*{Applications and Real-World Investigation}

\subsection*{Question 5 --- Hash functions in applications: passwords, sealed bids, Merkle trees}
\textbf{(a)} Password space $36^{8}\approx 2.82\cdot 10^{12}\approx 2^{41}$, $10^{6}$ users, verified in
\texttt{q5\_passwords.py} (Appendix~A.6).
\begin{enumerate}[label=(\roman*)]
\item \emph{Unsalted SHA-256:} searching the whole space once costs $2.82\cdot 10^{12}/10^{10}\approx
282$\,s $\approx \mathbf{4.7}$ \textbf{minutes}, and that single table breaks \emph{all} users at once,
because identical passwords hash identically and the table is precomputable (rainbow tables).
\item \emph{Per-user random salt, still SHA-256:} one chosen user still takes $\approx 4.7$ minutes (the
hash is just as fast), but all $10^{6}$ users now cost $10^{6}\times$ that $\approx \mathbf{8.9}$
\textbf{years}, because no table is shared across users.
\item \emph{Salted Argon2id} ($\approx 10^{4}$ guesses/s): one chosen user costs $2.82\cdot
10^{12}/10^{4}\approx \mathbf{8.9}$ \textbf{years} --- a $10^{6}\times$ slowdown versus the fast hash.
\end{enumerate}
A \textbf{salt} (stored, per-user, random) defeats precomputation and stops one effort from cracking
every user; a \textbf{pepper} (a site-wide secret kept outside the database, e.g.\ in an HSM) adds a
secret the attacker lacks even after stealing the database. \textbf{Memory-hard} functions force each
guess to use a large amount of RAM, which removes the massive parallelism of GPUs and ASICs (they have
little fast memory per core). Storing \emph{encrypted} passwords is worse than hashing, because
encryption is reversible: if the key leaks --- often from the same server --- every password is instantly
recovered, whereas a hash is one-way. Current guidance: OWASP recommends \textbf{Argon2id} with $\ge
19$\,MiB memory, iteration count $2$, parallelism $1$ (or bcrypt work factor $\ge 10$, or PBKDF2-HMAC-
SHA-256 with $\ge 600{,}000$ iterations) \cite{owasp}; NIST SP~800-63B requires a salt of at least 32
bits and storage with an approved key-derivation function, with an optional secret keyed hash
\cite{nist63b}. (Note that an 8-character alphanumeric space is only $2^{41}$, so even Argon2id merely
buys years for \emph{one} user --- password \emph{length} still matters most.)

\textbf{(b)} Sealed bids submitted as $h(\text{bid})$ for whole euro amounts $1$ to $10{,}000$. Chuck
simply computes $h(1),h(2),\dots,h(10{,}000)$ --- at most $\mathbf{10^{4}}$ hash computations --- and
matches Alice's published value to read her bid before the deadline. The commitment is not
\emph{hiding} because the input space is tiny and low-entropy (a trivial dictionary attack). \emph{Fix:}
commit to $h(\text{bid}\,\|\,r)$ with a long random nonce $r$, revealed only at opening. Then
\emph{hiding} comes from \textbf{preimage resistance} plus the high entropy of $r$ (the bid cannot be
guessed), and \emph{binding} comes from \textbf{collision/second-preimage resistance} (Alice cannot find
a second $(\text{bid}',r')$ with the same hash, so she cannot change her bid after committing).

\textbf{(c)} Merkle tree with $T(x,y)=(3x+5y+7)\bmod 101$, leaves $L_1{=}23,L_2{=}8,L_3{=}61,L_4{=}40$,
verified in \texttt{q5\_merkle.py} (Appendix~A.5).
\begin{enumerate}[label=(\roman*)]
\item $N_{12}=T(23,8)=116\bmod 101=15$; $N_{34}=T(61,40)=390\bmod 101=87$;
$R=T(15,87)=487\bmod 101=\boxed{83}$.
\item Inclusion proof for transaction 3: $L_3$ is the \emph{left} child of $N_{34}$, so the proof is
$[\,L_4=40\ (\text{right sibling}),\ N_{12}=15\ (\text{left uncle})\,]$. The light client recomputes
$T(N_{12},\,T(L_3,L_4)) = T(15,\,T(61,40)) = T(15,87)=83=R$ and accepts, having stored only the root.
\item If $L_3$ becomes $62$: $N_{34}'=T(62,40)=90$ and $R'=T(15,90)=\boxed{98}\neq 83$. The change is
detected because the recomputed root no longer matches the trusted root.
\item A block of $4096=2^{12}$ transactions needs an inclusion proof of $\log_2 4096 = \mathbf{12}$
hashes. The block header stores only the root (Stamp, Section~5.8.2), so a light client verifies any
transaction with a 12-hash proof instead of the whole block; and because each block header also commits
to the previous block's hash, altering an old transaction changes its root, hence its header hash, hence
every following block's ``previous hash'' --- invalidating the entire chain after it.
\end{enumerate}

\subsection*{Question 6 --- SHAttered: the first SHA-1 collision (2017)}
\textbf{(a) Event profile.}
\begin{table}[H]
\centering\small
\begin{tabularx}{\textwidth}{@{}P{3.6cm}L@{}}
\toprule
Date(s) & Announced 23 February 2017 (SHA-1 deprecated by NIST since 2011). \\
Who & Researchers at Google and CWI Amsterdam: M.~Stevens, E.~Bursztein, P.~Karpman, A.~Albertini, Y.~Markov \cite{shattered}. \\
Affected systems and scale & Everything relying on SHA-1 collision resistance: TLS/X.509 certificates (already being sunset), digital signatures, PGP, and content-addressed systems --- Git and Subversion (the WebKit SVN repository was corrupted by committing the two PDFs). SHA-1 was one of the most widely deployed hashes. \\
Hash and property & SHA-1 (160-bit, Merkle--Damg\aa rd); identical-prefix \emph{collision resistance}. \\
Identifier & shattered.io; CVE-2005-4900 (SHA-1 is not collision resistant). \\
Fix / outcome & Two PDFs with the same SHA-1 were published; migration to SHA-256 accelerated, Git shipped a collision-detecting SHA-1 (sha1dc) as the default in v2.13 (May 2017) \cite{git213}, and browsers already rejected SHA-1 certificates (Chrome~56, Jan 2017). \\
\bottomrule
\end{tabularx}
\end{table}

\textbf{(b) Technical root cause.} SHA-1's compression function was broken by \emph{differential}
cryptanalysis: the researchers built a two-block differential path (using disturbance vectors) so that a
carefully chosen difference between two inputs cancels out, producing the same chaining value. Combined
with enormous GPU computation this cost $\approx 2^{63.1}$ compression calls --- below the generic
$2^{80}$ --- and yielded two distinct PDF files with the \emph{same SHA-1 digest}. The property that
failed was \textbf{(identical-prefix) collision resistance}; preimage resistance was not broken. What the
attacker gains is two colliding documents: a signature or certificate on one is valid on the other, and
content-addressed stores that trust SHA-1 as an identifier can be poisoned. The \emph{retirement
timeline} was triggered step by step: Wang's 2005 theoretical attack ($2^{69}$) started deprecation;
SHAttered (2017, identical-prefix, $2^{63.1}$) proved it practical; and ``SHA-1 is a Shambles'' (2020,
\emph{chosen}-prefix, $2^{63.4}$, $\sim\$45$k) \cite{shambles} let attackers pick both prefixes freely
--- enough to forge realistic certificates and PGP identities --- which pushed NIST to announce full
retirement of SHA-1 by 2030 \cite{nistsha1}.

\textbf{(c) Classification.} The \textbf{hash function itself was broken by cryptanalysis}: the collision
exploited mathematical weaknesses in SHA-1's compression function, not a misuse by engineers. (This
contrasts with, say, the LinkedIn breach, where a sound idea was spoiled by fast, unsalted hashing.)
SHA-1 was used as intended --- for collision resistance --- and that property simply no longer held.

\textbf{(d) Response and lessons.} Vendors and standards bodies responded on several fronts: Git added
collision detection by default within months \cite{git213}, CAs and browsers had already scheduled the
SHA-1 certificate sunset (2016--2017), and NIST set a hard 2030 end date \cite{nistsha1}. The response
was effective but \emph{slow to start}: the community had warned since 2005, yet SHA-1 lingered for years,
and only the 2020 chosen-prefix result made migration urgent. A concrete control that would have reduced
the impact is \textbf{enforced crypto-agility with deprecation deadlines} --- standards and software that
make the hash a swappable parameter and bar weakened algorithms on a fixed schedule --- together with
collision detection enabled by default in content-addressed systems, so that a broken primitive can be
retired quickly rather than over a decade.

%=====================================================================
\section*{Design Decisions}

\subsection*{Question 7 --- Choosing hash-based mechanisms for UniServices}
\textbf{(a)} For 40{,}000 users over a 10-year lifetime:
\begin{table}[H]
\centering\small
\begin{tabularx}{\textwidth}{@{}P{2.2cm}P{4.3cm}LP{3.1cm}@{}}
\toprule
Requirement & Mechanism and parameters & Why not the naive alternative & One pitfall to avoid \\
\midrule
R1. Store passwords & Argon2id (19\,MiB, $t=2$, $p=1$), 128-bit per-user salt, optional pepper in an HSM \cite{owasp} & Plain/unsalted SHA-256 is fast and precomputable, so one GPU table breaks every account & Setting memory/iterations too low, or keeping the pepper on the same server as the hashes \\
R2. Authenticate API requests & HMAC-SHA-256 over the request, with a timestamp and nonce, keyed by the per-install shared secret & $h(K\|M)$ allows length-extension forgery; a plain hash authenticates nothing & No replay protection, or shipping the key in the app binary where it can be extracted \\
R3. Distribute exam-client updates & Code signing: ECDSA P-256 (or Ed25519) over SHA-256, public key pinned in the installed client & A SHA-256 checksum on the same server is rewritten by the same attacker who swaps the update & Leaving the signing key unprotected (sign in an HSM), or no rollback/version control \\
R4. Tamper-evident grade audit log & Hash chain / Merkle log: $H_i=\mathrm{SHA\text{-}256}(H_{i-1}\|\text{entry}_i)$, with roots periodically signed and published for the auditor & A plain log on the same server can be silently rewritten by an administrator & Not anchoring the head externally: an insider could recompute the whole chain unless roots are signed/published \\
\bottomrule
\end{tabularx}
\end{table}

\textbf{(b)} Post-quantum outlook for each choice:
\begin{itemize}
\item \textbf{R1 (Argon2id):} no change --- it is hash/symmetric-based, Grover only halves preimage
resistance, and 256-bit output keeps $\ge 128$-bit security.
\item \textbf{R2 (HMAC-SHA-256):} no change --- HMAC is symmetric and remains secure against quantum
attackers at these sizes.
\item \textbf{R3 (ECDSA code signing):} \textbf{must change} --- Shor's algorithm breaks ECDSA, and a
10-year client with a long-lived trust anchor makes migration urgent; move to a hash-based or
lattice signature (SLH-DSA/LMS or ML-DSA).
\item \textbf{R4 (hash chain):} the hashing needs no change (SHA-256 is fine), but the \emph{signature}
on the published roots must migrate to a post-quantum scheme so the auditor's verification stays sound.
\end{itemize}

\section*{Conclusion}
Hash functions underpin passwords, downloads, MACs, signatures and blockchains, and their security rests
on three properties plus the birthday bound, which forces a hash to be twice as long as a symmetric key.
The toy checksum showed why linearity and short output make a function trivially invertible, and the
survey showed the real split in the field: MD5 and SHA-1 are broken by cryptanalysis and retired, while
SHA-256, SHA-512 and SHA3-256 remain intact. The constructions matter as much as the function:
Merkle--Damg\aa rd exposes its whole state and so allows length extension, which HMAC and the sponge
avoid, and the SHAttered case confirmed that a mathematically broken primitive endangers every signature
built on it. The applications reinforced that correct \emph{use} is just as important as a strong
function --- salting and memory-hard hashing for passwords, high-entropy commitments for sealed bids, and
signed Merkle roots for tamper-evidence. Finally, the post-quantum view is reassuring for hashing:
256-bit hashes survive Grover and BHT, and hash-based signatures (SLH-DSA, LMS/XMSS) inherit that
strength, so the real migration work for a long-lived system lies in its signatures, not its hashes.

%=====================================================================
\clearpage
\section*{Appendix A --- Python Verification: Chosen Values and Outputs}
The six scripts are in the repository \url{https://github.com/Diellart-Mulolli/hash-lab} (see its README
for the running guide). They were run with Python~3.12.3 (standard library only). The outputs below are
copied unchanged from the terminal, and the key results are cross-checked independently: the custom
SHA-256 in A.3 is validated against Python's \texttt{hashlib}, and the HMAC structure in A.4 against
Python's \texttt{hmac} module.

\subsection*{A.1 \texttt{q1\_toy\_hash.py} (Question 1(b))}
Chosen values: toy hash $h(X)=(3X_0+2X_1+X_2)\bmod 256$; message $\mathrm{BID}=(66,73,68)$; targeted
prefix \texttt{OK}$=(79,75)$. Command: \texttt{python3 q1\_toy\_hash.py}
\begin{Verbatim}[fontsize=\footnotesize, frame=single]
message BID = (66, 73, 68) = b'BID'
(i)   h(BID) = (3*66 + 2*73 + 68) mod 256 = 412 mod 256 = 156

(ii)  second preimage: (67, 73, 65) = b'CIA', h = 156  (equal: True)
      total 3-byte messages with this hash: 65536 (= 256^2, one for every (X0,X1))

(iii) targeted preimage starting 'OK': X2 = (156 - 3*79 - 2*75) mod 256 = 25
      message = (79, 75, 25) = b'OK\x19', h = 156  (equal: True)
\end{Verbatim}
\textbf{Check.} $h(\mathrm{BID})=156$ matches the hand computation; the second preimage \texttt{CIA} and
the targeted preimage both hash to $156$; and a full scan confirms $65536$ messages share the value,
showing the formula is many-to-one and linear.

\subsection*{A.2 \texttt{q2\_birthday.py} (Questions 1(c), 2(b), 2(c))}
Chosen values: birthday formula $P(k)=1-e^{-k(k-1)/2^{n+1}}$; $n=8$ with $k=16$; attacker rate
$10^{12}$/s; exponents $64,80,128$; 32 either/or phrasings. Command: \texttt{python3 q2\_birthday.py}
\begin{Verbatim}[fontsize=\footnotesize, frame=single]
Q1(c): 8-bit toy hash
  P(16 messages) = 1 - exp(-16*15 / 2^9) = 1 - exp(-240/512) = 0.3742  (~37.4%)
  k for P~=50%: 1.18 * 2^(8/2) = 1.18 * 16 = 18.9  (~19 messages)
  smallest k with P>=0.5 (exact from the formula): 20

Q2(b): attacker at 1e12 hashes/second
  2^64: 1.84e+07 s = 214 days
  2^80: 1.21e+12 s = 3.83e+04 years
  2^128: 3.4e+26 s = 1.08e+19 years
  for comparison, 2^80 / 1e12 = 3.83e+04 years
  SHAttered (2017) used ~2^63.1 compression calls; 2^63.1/1e12 = 114.41 days at this imagined rate

Q2(c): birthday attack on signatures
  32 independent either/or phrasings -> 2^32 = 4294967296 variants of each message
  collision between the 'good' and 'evil' sets needs ~2^(n/2) variants;
  2^32 variants defeat any hash with n <= 64 bits (2^(64/2) = 2^32).
\end{Verbatim}
\textbf{Check.} $P(16)\approx 37.4\%$ and the $50\%$ threshold near $k\approx 19$ (rule of thumb) / $20$
(exact) match Question~1(c); the $2^{64},2^{80},2^{128}$ times and the $2^{32}$ variant count match
Questions~2(b)--(c).

\subsection*{A.3 \texttt{q3\_length\_extension.py} (Question 3(b))}
Chosen values: 16-byte key \texttt{"sixteen-byte-key"}; $M=\texttt{"amount=10\&to=alice"}$; appended
suffix \texttt{"\&to=trudy"}. The script implements SHA-256 from scratch so it can resume from a captured
digest. Command: \texttt{python3 q3\_length\_extension.py}
\begin{Verbatim}[fontsize=\footnotesize, frame=single]
sanity: custom SHA-256 matches hashlib: True
intercepted M   = b'amount=10&to=alice' (18 bytes)
intercepted tag = 2f54ff5373411d83d51cc5233569996c17d1e97323567152ff07ef6450cb449a

glue padding = 0x80 then zeros then 00..0110 (30 bytes; 0x110 = 272 bits)
forged M' = M || glue(30 B) || b'&to=trudy'  (total 57 bytes)
forged tag (no key)   = 6abffc5fdeeda4397a8d76ba769e43000aa94e933bee39ce881e3d7a7723c1ac
server tag (with key) = 6abffc5fdeeda4397a8d76ba769e43000aa94e933bee39ce881e3d7a7723c1ac
forgery accepted: True
server uses the LAST to= value, so the request now pays to=trudy.
\end{Verbatim}
\textbf{Check.} The from-scratch SHA-256 matches \texttt{hashlib} (first line), and the forged tag
computed \emph{without} the key equals the tag the server computes \emph{with} the key --- a successful
length-extension forgery.

\subsection*{A.4 \texttt{q4\_hmac.py} (Question 4(b))}
Chosen values: $K=\texttt{key}=\texttt{0x6B\,0x65\,0x79}$; $B=64$; $ipad=\texttt{0x36}$,
$opad=\texttt{0x5C}$; 1000-byte message for the block count. Command: \texttt{python3 q4\_hmac.py}
\begin{Verbatim}[fontsize=\footnotesize, frame=single]
key K = b'key' = 0x6B 0x65 0x79

(i)   K is shorter than 64 bytes, so it is right-padded with zero bytes to 64 bytes.

(ii)  first three bytes (the rest are the pad byte itself, since 0x00 XOR pad = pad):
  K XOR ipad: 0x5D 0x53 0x4F ... then 0x36 x61
  K XOR opad: 0x37 0x39 0x25 ... then 0x5C x61
    ipad: k=0x6B^0x36=0x5D ; e=0x65^0x36=0x53 ; y=0x79^0x36=0x4F
    opad: k=0x6B^0x5C=0x37 ; e=0x65^0x5C=0x39 ; y=0x79^0x5C=0x25

(iii) HMAC(M,K) = SHA256( (K XOR opad) || SHA256( (K XOR ipad) || M ) )

(iv)  SHA-256 of a 1000-byte message: ceil((1000+9)/64) = 16 blocks
      HMAC inner SHA-256 over (K^ipad)||M: ceil((1064+9)/64) = 17 blocks
      HMAC outer SHA-256 over (K^opad)||digest: ceil((96+9)/64) = 2 blocks
      HMAC total = 17 + 2 = 19 blocks vs 16 for a plain hash (+3 blocks, ~2 extra compressions, negligible).

cross-check on M=b'transfer 100 EUR to Bob': manual == hmac.new: True
  HMAC = ca9c97583730acf2db3b30015cfcdb15b70ae3aace136d7ead2d43050b95d2e3
\end{Verbatim}
\textbf{Check.} The padded key length, the first bytes of $K\oplus ipad$/$K\oplus opad$, and the block
counts match Question~4(b); the manual two-pass construction equals Python's \texttt{hmac.new}.

\subsection*{A.5 \texttt{q5\_merkle.py} (Question 5(c))}
Chosen values: combiner $T(x,y)=(3x+5y+7)\bmod 101$; leaves $(23,8,61,40)$; altered leaf $61\to 62$.
Command: \texttt{python3 q5\_merkle.py}
\begin{Verbatim}[fontsize=\footnotesize, frame=single]
(i) building the tree
  N12 = T(23, 8)  = (3*23 + 5*8  + 7) mod 101 = 116 mod 101 = 15
  N34 = T(61, 40) = (3*61 + 5*40 + 7) mod 101 = 390 mod 101 = 87
  R   = T(N12,N34)= (3*15 + 5*87 + 7) mod 101 = 487 mod 101 = 83

(ii) inclusion proof for transaction 3 (leaf L3, which is the LEFT child of N34)
  proof = [ L4 = 40 (sibling, right), N12 = 15 (uncle, left) ]
  light client: T(N12, T(L3, L4)) = T(15, T(61,40)) = T(15,87) = 83  (matches root 83: True)

(iii) transaction 3 altered: L3 = 61 -> 62
  N34' = T(62, 40) = 90,  R' = T(15, 90) = 98
  root changed 83 -> 98: detected because the stored/signed root no longer matches (True)

(iv) proof size for a block of 4096 transactions
  depth = log2(4096) = 12, so an inclusion proof has 12 hashes
\end{Verbatim}
\textbf{Check.} Root $R=83$; the inclusion proof for transaction~3 recomputes to $83$; the altered leaf
changes the root to $98$; and a 4096-leaf tree needs 12 proof hashes --- matching Question~5(c).

\subsection*{A.6 \texttt{q5\_passwords.py} (Question 5(a))}
Chosen values: charset $36$, length $8$ ($36^8$ space); $10^6$ users; $10^{10}$/s fast hash,
$10^{4}$/s Argon2id. Command: \texttt{python3 q5\_passwords.py}
\begin{Verbatim}[fontsize=\footnotesize, frame=single]
password space = 36^8 = 2,821,109,907,456 ~= 2^41
users = 1,000,000

(i) unsalted SHA-256:
    search the whole space once = 2,821,109,907,456/1e+10 = 0.0784 hours
    the same work breaks ALL users (identical passwords hash identically; precomputable).

(ii) per-user random salt, still SHA-256:
    one chosen user  = 0.0784 hours (same as (i): the hash is still fast)
    all 1,000,000 users = 1,000,000 x that = 8.94 years (no shared table any more)

(iii) salted Argon2id (~1e4 guesses/s):
    one chosen user = 2,821,109,907,456/1e+04 = 8.94 years
    slowdown vs fast hash = 1e+06x
\end{Verbatim}
\textbf{Check.} Unsalted SHA-256 falls in minutes for all users; salting makes the all-users cost
$\approx 8.9$ years; Argon2id makes even one user $\approx 8.9$ years --- the figures used in
Question~5(a).

\clearpage
\begin{thebibliography}{99}
\footnotesize\setlength{\itemsep}{1pt}\setlength{\parsep}{0pt}\setlength{\parskip}{0pt}
\bibitem{stamp}
M.~Stamp, \emph{Information Security: Principles and Practice}, 3rd~ed. Hoboken, NJ, USA: Wiley, 2022, ch.~5.

\bibitem{shattered}
M.~Stevens, E.~Bursztein, P.~Karpman, A.~Albertini, and Y.~Markov, ``The first collision for full SHA-1,'' in \emph{Advances in Cryptology -- CRYPTO 2017} (LNCS 10401). Cham, Switzerland: Springer, 2017, pp.~570--596. [Online]. Available: \url{https://shattered.io}

\bibitem{shambles}
G.~Leurent and T.~Peyrin, ``SHA-1 is a shambles: First chosen-prefix collision on SHA-1 and application to the PGP web of trust,'' in \emph{Proc.\ 29th USENIX Security Symp.}, 2020, pp.~1839--1856.

\bibitem{stevenscp}
M.~Stevens, A.~Lenstra, and B.~de~Weger, ``Chosen-prefix collisions for MD5 and applications,'' \emph{Int.\ J.\ Appl.\ Cryptogr.}, vol.~2, no.~4, pp.~322--359, 2012.

\bibitem{sha2coll}
Y.~Li, F.~Liu, and G.~Wang, ``New records in collision attacks on SHA-2,'' in \emph{Advances in Cryptology -- EUROCRYPT 2024} (LNCS 14651). Cham, Switzerland: Springer, 2024, pp.~158--186, Cryptology ePrint Archive, Paper 2024/349. [Online]. Available: \url{https://eprint.iacr.org/2024/349}

\bibitem{sha3coll}
J.~Guo, G.~Liao, G.~Liu, M.~Liu, K.~Qiao, and L.~Song, ``Practical collision attacks against round-reduced SHA-3,'' \emph{J.\ Cryptology}, vol.~33, no.~1, pp.~228--270, 2020, doi: \href{https://doi.org/10.1007/s00145-019-09313-3}{10.1007/s00145-019-09313-3}.

\bibitem{nistsha1}
NIST, ``NIST retires SHA-1 cryptographic algorithm,'' News, 15 Dec.~2022. [Online]. Available: \url{https://www.nist.gov/news-events/news/2022/12/nist-retires-sha-1-cryptographic-algorithm} (accessed Oct.~20, 2026).

\bibitem{git213}
GitHub, ``Git 2.13 has been released,'' The GitHub Blog, 10 May 2017. [Online]. Available: \url{https://github.blog/open-source/git/git-2-13-has-been-released/} (accessed Oct.~20, 2026).

\bibitem{fips205}
\emph{Stateless Hash-Based Digital Signature Standard}, FIPS PUB~205, Nat.\ Inst.\ Standards Technol., Gaithersburg, MD, USA, Aug.~2024, doi: \href{https://doi.org/10.6028/NIST.FIPS.205}{10.6028/NIST.FIPS.205}.

\bibitem{sp800208}
D.~Cooper \emph{et al.}, ``Recommendation for stateful hash-based signature schemes,'' Nat.\ Inst.\ Standards Technol., Gaithersburg, MD, USA, NIST SP~800-208, Oct.~2020, doi: \href{https://doi.org/10.6028/NIST.SP.800-208}{10.6028/NIST.SP.800-208}.

\bibitem{owasp}
OWASP, ``Password Storage Cheat Sheet,'' OWASP Cheat Sheet Series, 2024. [Online]. Available: \url{https://cheatsheetseries.owasp.org/cheatsheets/Password_Storage_Cheat_Sheet.html} (accessed Oct.~20, 2026).

\bibitem{nist63b}
D.~Temoshok \emph{et al.}, ``Digital identity guidelines: Authentication and authenticator management,'' Nat.\ Inst.\ Standards Technol., Gaithersburg, MD, USA, NIST SP~800-63B, 2025, doi: \href{https://doi.org/10.6028/NIST.SP.800-63b}{10.6028/NIST.SP.800-63b}.

\bibitem{rfc2104}
H.~Krawczyk, M.~Bellare, and R.~Canetti, ``HMAC: Keyed-hashing for message authentication,'' Internet Eng.\ Task Force, RFC~2104, Feb.~1997. [Online]. Available: \url{https://www.rfc-editor.org/rfc/rfc2104}

\bibitem{fips202}
\emph{SHA-3 Standard: Permutation-Based Hash and Extendable-Output Functions}, FIPS PUB~202, Nat.\ Inst.\ Standards Technol., Gaithersburg, MD, USA, Aug.~2015, doi: \href{https://doi.org/10.6028/NIST.FIPS.202}{10.6028/NIST.FIPS.202}.
\end{thebibliography}

\end{document}
